\section{Discussion}
\label{sec:discussion}

\subsection{Why the repair works where it does}
\label{sec:discussion-why}

The wrapped models gain coverage in the stress windows because the stress windows are the case the update was built for: a regime in which the model understates the tail by a scale factor that persists for months. Historical simulation is the clearest example. Its window of 250 days still holds the calm year when the crisis begins, so its empirical quantile is a scale error of a factor that grows with the ratio of the crisis volatility to the calm volatility; the working level falls until the quantile of that old window reaches the new losses. The scale-sweep experiment says how far this can go: a scale error of 1.5 at the 1\% level is repaired in direction but not fully, because the fixed point lies close enough to the lower edge of the band that the spread of the level makes the band bind, and the identity of \cref{prop:identity} accounts for the excess of the hit rate to within 0.01 points. The models whose own filter follows the volatility (GARCH-t, FHS, the extreme-value model) have a smaller error to repair, gain fewer passes per cell, and at the 1\% level pay a small score cost, within the tolerance, a pattern consistent with a level that fluctuates around a fixed point that is already near the target.

That the independence test does not fail after the wrap is not a property of the method. \Cref{prop:independence} shows that it can fail with certainty on a sequence that the identity allows, and the stress windows contain few enough hits (between 2 and 17 per cell for the wrapped models of the main table) that the independence test has little power to reject. The result of the experiment is that the wrapper did not cause a rejection in \tenIndRunsAll{} runs, and that the unwrapped models had \tenIndFailUnwAll{} rejections of independence; neither number says that the wrapper repairs clustering. A smaller number of hits gives the independence statistic fewer pairs of consecutive hits to count, so a lower rate of hits may by itself lower the statistic without repairing any dependence.

\subsection{The failure cases}
\label{sec:discussion-failures}

Four failures are in the results and should be read as properties of the method. The delta-normal model gains no stress window at the 1\% level; it shares with the other models a hit rate that improves, and differs by a tail so thin that the working level sits at the band edge on \clipDnRollingOneEqCrisis{} days of the crisis window and the hits that remain fail the tests. The historical-simulation models still fail the Kupiec test at the 1\% level in the crisis window and at both levels in the pandemic window: a window of 104 days that contains a crash, with a level that needs about 100 quiet days to relax after one hit, is a case where the recalibration is late by construction. The wrapper does not repair the ES (the mean $Z_2$ remains negative), which is the consequence of paying no attention to the ES in the update and of the fact that the tails of the models are not Pareto (\cref{prop:es}). And for the models that already track volatility, the generator experiments show a loss of accuracy against the true VaR in \genErrWorseCells{} of \genErrWorseDen{} cells: the update trades some of the accuracy of a good estimate for the coverage of a bad one, and the supervisor's tests, which are on coverage, would not show the trade.

\subsection{What the comparison with other recalibrations implies}
\label{sec:discussion-comparison}

The comparison in \cref{sec:results-comparators} removes an easy conclusion. If the wrapper were the only way to move the coverage of a model towards its target, its gain in cells passed would be a property of the method; it is not, since a rolling split-conformal scale passes as many cells and a quantile tracker at a larger step passes more. The reason is that the repair of coverage under a persistent scale error is the common content of all these updates, which share the feature that a hit moves the working quantity one way and a quiet day the other. The differences between them are second-order: the rolling scale has no memory of the direction of its error and moves with the sample quantile of the last 250 scores, whose own noise at the 1\% level (the second largest of 250 scores, and the sixth largest at 2.5\%) moves the VaR by large factors, a plausible reason for its \cmpRollTenFzOver{} cells over the score tolerance, the worst of them in the short pandemic window where a crash lifts the scale to \cmpRollWorstScaleMax{}; the tracker has a step with the dimension of a score and a pass count that rises with it; the level wrapper keeps the VaR and the ES on one distribution and pays for the speed of its step in the score, as the ablation shows. What the wrapper offers a bank is not a larger number of passed windows but a recalibration that leaves the model's own distribution in charge of the ES, within a score tolerance, with a step that has a meaning in probability (a hit moves the level by $\gamma(1-\alpha)$). Whether that is worth a model of its own is a decision for the validator, and the data here do not make it.

\subsection{What a pass count measures}
\label{sec:discussion-passcount}

A bank reads a back-test as a binary outcome: the model passes or it does not. \Cref{sec:results-dispersion} shows that the outcome depends on the dispersion of the hit count as much as on its mean. The update ties the count of hits to its target through the identity of \cref{prop:identity}, whatever the model, and a count tied in this way is accepted by tests that assume independent hits. A validator who runs the Kupiec or the conditional coverage test on a recalibrated VaR therefore tests a null that the recalibration has changed, and the run passes more often than a correct unrecalibrated model would. The consequence is not specific to this wrapper: the rolling scale and the quantile tracker show a weaker under-dispersion in the generators (ratios \dispRatioRollMin{} to \dispRatioRollMax{} and \dispRatioTrkMin{} to \dispRatioTrkMax{}, against \dispRatioAciMin{} to \dispRatioAciMax{} for the wrapper), and a recalibration that feeds the hit back into the forecast can be expected to show it. Three readings follow. The hit rate and its distance from the target, which the tests do not enter, are the evidence of a repair that survives. The pass counts of a recalibrated model are to be read against a null simulated with the same recalibration, as in the recount of \cref{sec:results-dispersion}, where the gain falls from \tenPassAci{} against \tenPassUnw{} cells to \dispCWrapCal{} against \dispCUnwExact{}. And the supervisor's traffic light reads a count of exceptions and would, by the same argument, place a count held near its target in the green zone more often than a count of independent hits (the experiments tabulate hits, not zones, so this is not measured here), which is a reason to ask the bank for the recalibration rule and not only for the count.

\subsection{Implications for the validation of a bank's model}
\label{sec:discussion-implications}

A validator who finds that a model fails in a stress window has two instruments in hand after this work. The first is to read the wrapper as a diagnostic: the path of the working level, the days at the band edge and the mean level tell how far the model is from its target in scale. A model whose level sits at the lower edge for months, as the delta-normal model does, has an error that no step repairs; a model whose level moves by a few tenths of the target and returns has a localised failure. The second is the repair itself, for the models whose structure cannot be changed quickly. The repair does not substitute for the capital multiplier of the standard \citep{bcbs1996}, which rises with the count of exceedances over 250 days and is a consequence applied after the fact; it changes the VaR that produces the exceedances. A bank that adopts it would present to the supervisor the unwrapped model's back-test, the wrapped model's back-test, and the level path, since the supervisor reads the VaR that the bank reports and the reported VaR would be the wrapped one. The wrapper's dependence on realised losses of the day also means that a bank has to record the hit before the next day's forecast; the hit of a late-booked day is out of scope of the method and of its test.
